\section{Every labelled-target predecessor}\label{sec:inverse}

The inverse problem fixes the whole target function, including its
endpoint labels. Its solution is not a test involving an unknown source
map: the following construction lists all sources directly from the
target.

For a nonnegative integer $d$, define the finite Laurent polynomial
\begin{equation}\label{eq:gap_poly}
 P_d(u)=\sum_{\substack{a,b\ge0\\a+b\le d}}
             \binom{d}{a+b}u^{b-a}.
\end{equation}
The notation $[u^j]Q(u)$ denotes the coefficient of $u^j$ in a Laurent
polynomial $Q$.

\begin{theorem}\label{thm:inverse}
A target $g=f_{Z,W}$ failing~\eqref{eq:image} has no predecessor.
Otherwise put $z_0=0$ and
\[
 D_i=\{z_{i-1}+1,\ldots,w_i-1\},\qquad
 d_i=|D_i|=w_i-z_{i-1}-1\quad(1\le i\le k).
\]
Every predecessor is obtained uniquely as follows.
\begin{enumerate}
\item Start with $X=Z$ and $A=W$.
\item In each $D_i$, select disjoint sets of extra $X$-sites and extra
      $A$-sites such that every selected $X$-site precedes every selected
      $A$-site. Sites not selected belong to neither set. Add no other
      sites outside these free gaps.
\item Retain the choices with $\delta=|A|-|X|\in\{0,1\}$. Set
      $Y=A$ if $\delta=0$ and $Y=A\setminus\{n\}$ if $\delta=1$.
      The reconstructed source is $f_{Y,X}$.
\end{enumerate}
In particular,
\begin{equation}\label{eq:fibre}
 |T^{-1}(g)|=\bigl([u^0]+[u^1]\bigr)
                    \prod_{i=1}^{k}P_{d_i}(u).
\end{equation}
\end{theorem}

\begin{proof}
The zero-fibre assertion is Lemma~\ref{lem:product}. Suppose the image
condition holds and $T(f_{Y,X})=g$, with $A=Y\cup\{n\}$.
The product calculation forces $z_i\in X$ and $w_i\in A$.
For every strict pair $w_i<z_i$, it also forces
\begin{equation}\label{eq:forbidden}
 X\cap[w_i,z_i)=\varnothing,\qquad
 A\cap(w_i,z_i]=\varnothing.
\end{equation}
The first condition is necessary for $e_X(w_i)=z_i$. For the second,
an extra $a\in A$ with $w_i<a\le z_i$ would also have ceiling
$z_i$, contradicting that $w_i$ is the final domain endpoint of that
output block. These prohibitions, the forced endpoints and the $D_i$
partition all sites of $[n]$.

A free site cannot belong to both supports, because the identity
$X\cap A=W\cap Z$ would give a further target anchor. In $D_i$,
no selected $A$-site may precede a selected $X$-site. Indeed, if
$a<x$ with $a\in A$, $x\in X$ and both in $D_i$, then
\[
 z_{i-1}<a\le e_X(a)\le x<w_i\le z_i.
\]
The extra value $e_X(a)$ lies strictly between two consecutive target
values, a contradiction. Thus the ordered-colour condition is necessary.

It is also sufficient. Make any choices in the stated free gaps and
include the forced endpoints. Every selected $A$-site in
$(z_{i-1},w_i]$ lies after all extra $X$-sites in that gap and has
ceiling exactly $z_i$ in $X$. Its last such site is $w_i$.
The product $e_X\circ e_A$ therefore has precisely the prescribed
values $z_i$ and kernel endpoints $w_i$, and equals $g$.

It remains to impose the source-rank condition. The forced sets $Z,W$
have the same size, even when they meet at an anchor. If $a_i,b_i$
extra sites of the two respective colours are chosen in $D_i$, then
\[
 |A|-|X|=\sum_{i=1}^{k}(b_i-a_i).
\]
We now use the published completed-support reconstruction
\cite[Section 5.1, after Lemma 5.12]{stein2025order}: the allowed
differences are zero and one, recovering $Y=A$ and
$Y=A\setminus\{n\}$, respectively. This support-rank fact is not a
new part of the inverse theorem. In the one branch, $|A|=|X|+1\ge2$,
so removing $n$ does not produce an empty image. Each retained pair
gives exactly one source, and the necessity argument recovers its unique
gap description.

For fixed $a_i,b_i$, choose $a_i+b_i$ positions of $D_i$; their first
$a_i$ positions must have colour $X$ and their last $b_i$ colour $A$.
There are exactly $\binom{d_i}{a_i+b_i}$ choices. The Laurent weight
records their contribution $b_i-a_i$ to the rank difference. Independent
gap choices multiply, and selecting the zero and one coefficients gives
\eqref{eq:fibre}. This count and the full decoder use no temporal theorem.
\end{proof}

Empty gaps cause no exception: $P_0(u)=1$ and contribute no extra
sites. At $n=1$, the sole target has $d_1=0$, so the formula gives its
single predecessor. Formula~\eqref{eq:fibre} is a finite target-specific
evaluation; no maximization over targets is asserted.
